\documentclass[12pt, a4paper]{article}

\usepackage{amsmath, amssymb, amsthm}
\usepackage{geometry}
\usepackage{hyperref}
\usepackage{titlesec}
\usepackage{parskip}
\usepackage{microtype}
\usepackage{lmodern}
\usepackage{xcolor}
\usepackage{booktabs}
\usepackage{enumitem}
\usepackage{listings}
\usepackage{array}

\geometry{margin=1.2in}

\definecolor{deepblue}{RGB}{20, 40, 100}
\definecolor{muted}{RGB}{90, 90, 90}
\definecolor{codegray}{RGB}{245, 245, 245}
\definecolor{codegreen}{RGB}{0, 120, 0}
\definecolor{codepurple}{RGB}{120, 0, 200}

\hypersetup{
  colorlinks=true,
  linkcolor=deepblue,
  urlcolor=deepblue,
  citecolor=deepblue
}

\titleformat{\section}{\large\bfseries\color{deepblue}}{{\thesection}}{1em}{}
\titleformat{\subsection}{\normalsize\bfseries\color{deepblue}}{{\thesubsection}}{1em}{}
\titleformat{\subsubsection}{\normalsize\itshape\color{deepblue}}{{\thesubsubsection}}{1em}{}

\lstdefinestyle{python}{
  backgroundcolor=\color{codegray},
  basicstyle=\ttfamily\small,
  breaklines=true,
  commentstyle=\color{codegreen},
  keywordstyle=\color{deepblue}\bfseries,
  stringstyle=\color{codepurple},
  language=Python,
  frame=single,
  framesep=4pt,
  xleftmargin=10pt,
  xrightmargin=10pt,
  numbers=left,
  numberstyle=\tiny\color{muted},
  numbersep=8pt,
  showspaces=false,
  tabsize=4,
}

\theoremstyle{definition}
\newtheorem{definition}{Definition}
\newtheorem{decision}{Design Decision}

\newcommand{\R}{\mathbb{R}}
\newcommand{\Z}{\mathbb{Z}}
\newcommand{\Action}{\mathcal{A}}
\newcommand{\Fix}{\mathrm{Fix}}
\newcommand{\bfq}{\mathbf{q}}
\newcommand{\bfa}{\mathbf{a}}
\newcommand{\bfb}{\mathbf{b}}
\newcommand{\SO}{\mathrm{SO}}
\newcommand{\eps}{\varepsilon}


\begin{document}

% ---------- TITLE ----------
\begin{center}
  {\LARGE \textbf{Three-Dimensional Choreographic Orbits:\\[4pt]
  Computational Specification}}\\[1.6em]
  {\large \textcolor{muted}{CompSpec — v0.2 \quad|\quad Companion to RDD v0.2 and MathFrame v0.1}}\\[0.6em]
  {\normalsize \textcolor{muted}{e5bed97b3349 \quad\textbullet\quad March 2026}}
\end{center}

\vspace{1em}
\noindent\rule{\textwidth}{0.4pt}
\vspace{0.5em}

\tableofcontents

\vspace{0.5em}
\noindent\rule{\textwidth}{0.4pt}

% ---------- 1. SOFTWARE STACK ----------
\section{Software Stack and Hardware Target}

\subsection{Primary framework: JAX}

\begin{decision}
All action evaluation and gradient computation will be implemented in
\textbf{JAX} \texttt{(jax>=0.4.25)}, rather than NumPy, PyTorch, or SciPy.
\end{decision}

JAX is the correct choice for three compounding reasons. First, the action
functional $\Action[\gamma]$ is a smooth function of the Fourier coefficients;
JAX's \texttt{jax.grad} and \texttt{jax.hessian} provide exact gradients and
Hessians at machine precision without finite-differencing. Second,
\texttt{jax.vmap} vectorises the action evaluation over an entire batch of
candidate coefficient vectors with zero loop overhead, enabling thousands of
random initialisations to be evaluated simultaneously. Third, \texttt{jax.jit}
compiles the entire pipeline — Fourier synthesis, pairwise distance computation,
action summation — to XLA and dispatches it to the GPU without data movement
overhead.

\subsection{Supporting libraries}

\begin{center}
\begin{tabular}{lll}
\toprule
\textbf{Library} & \textbf{Version} & \textbf{Role} \\
\midrule
\texttt{jax[cuda12]} & $\geq$\,0.4.25 & Core autodiff, JIT, GPU dispatch \\
\texttt{optax}       & $\geq$\,0.2.2  & L-BFGS and Adam optimisers \\
\texttt{scipy}       & $\geq$\,1.13   & Newton--Raphson, eigenvalues \\
\texttt{snappy}      & $\geq$\,3.1    & Knot invariants (Alexander, HOMFLY-PT) \\
\texttt{pyvista}     & $\geq$\,0.44   & 3-D tube renders, symmetry-axis overlays \\
\texttt{ffmpeg-python} & $\geq$\,0.2  & 4K/60fps video export \\
\texttt{numpy}       & $\geq$\,1.26   & Data I/O, catalogue serialisation \\
\texttt{h5py}        & $\geq$\,3.11   & HDF5 storage for orbit catalogue \\
\bottomrule
\end{tabular}
\end{center}

\subsection{Hardware target}

The primary compute target is a single workstation equipped with an NVIDIA
RTX\,5090 (Blackwell, 32\,GB VRAM) and 192\,GB system RAM.  Key
implications:

\begin{itemize}[leftmargin=1.8em]
  \item The Fourier truncation limit $K$ can be pushed to $K = 256$ or beyond
        without memory pressure, capturing highly irregular curve geometries
        that lower-$K$ truncations would miss.
  \item With \texttt{vmap}, a batch of $B = 8192$ random initialisations can be
        evaluated in a single forward pass; the GPU's 3352 tensor cores
        parallelise all $B$ action evaluations concurrently.
  \item The 192\,GB system RAM allows the full symmetry-constraint manifolds for
        larger $N$ to be pre-enumerated and held in RAM during search.
\end{itemize}

% ---------- 2. MATHEMATICAL ENCODING ----------
\section{Mathematical Encoding}

\subsection{Fourier representation of the curve}

The single curve $\gamma : [0,T] \to \R^3$ is represented by a truncated
real Fourier series:
\[
  \gamma(t) = \bfa_0
  + \sum_{k=1}^{K} \!\left[
      \bfa_k \cos\!\left(\tfrac{2\pi k t}{T}\right)
    + \bfb_k \sin\!\left(\tfrac{2\pi k t}{T}\right)
  \right],
  \qquad \bfa_k, \bfb_k \in \R^3.
\]
The parameter vector is $\theta = (\bfa_0, \bfa_1, \bfb_1, \ldots,
\bfa_K, \bfb_K) \in \R^{3(2K+1)}$.  The centre-of-mass condition
$\sum_i \bfq_i = 0$ for all $t$ enforces $\bfa_0 = \mathbf{0}$, reducing
$\theta \in \R^{6K}$.  The time discretisation uses $M = 8K$ uniformly spaced
quadrature points; the action integral is approximated by the rectangle rule,
whose error is $O(K^{-p})$ for a curve of Sobolev regularity $p$.

\begin{decision}
Default truncation: $K = 64$ for Phase\,1 searches.  Raise to $K = 128$
for refinement and $K = 256$ for near-collision orbits flagged during Phase\,1.
Monitor convergence via $\|\bfa_K\|, \|\bfb_K\| < 10^{-8}$.
\end{decision}

\subsection{Symmetry constraint encoding}

Let $G \subset O(3) \times \Z/N\Z$ be the target symmetry group.  Each
generator $(R, \sigma) \in G$ acts on $\theta$ by:
\[
  (R,\sigma) \cdot \gamma(t) = R\,\gamma\!\left(t - \sigma\,\tfrac{T}{N}\right).
\]
In Fourier space this becomes a \emph{linear} constraint on $\theta$.
Concretely, applying $R$ to the coefficient vectors and shifting $t$ by
$\sigma T/N$ multiplies mode $k$ by the phase factor $e^{2\pi i k \sigma/N}$.
The fixed-point subspace $\Fix(G)$ is the intersection of the null spaces of
$(I - \rho(g))$ for each generator $g \in G$, where $\rho$ is the induced
representation on $\R^{6K}$.

\subsubsection{Worked example: $C_{3v}$ with $N = 3$}

The group $C_{3v}$ is generated by a $120^\circ$ rotation $R_3$ about the
$z$-axis and a reflection $\sigma_v$ through the $xz$-plane.  With $N = 3$
and phase shift $\sigma = 1$:

\begin{enumerate}[leftmargin=1.8em]
  \item $R_3 \cdot \gamma(t) = R_3\,\gamma(t - T/3)$: in the helicity basis
        $z_\pm = x \pm iy$, mode $k$ acquires phase $e^{2\pi i(1-k)/3}$ in the
        $xy$-plane and $e^{-2\pi ik/3}$ on the $z$-axis.  Survival requires
        these phases to equal 1, giving $k \equiv 1 \pmod{3}$ for
        $xy$-components and $k \equiv 0 \pmod{3}$ for $z$-components
        (MathFrame Proposition~4.1).
  \item $\sigma_v \cdot \gamma(t) = \sigma_v \gamma(t)$: reflects $y \to -y$,
        forcing $(\bfa_k)_y = 0$ and $(\bfb_k)_x = 0$ for all surviving modes.
\end{enumerate}

The net effect is that $\dim\Fix(C_{3v}) = 4K/3 + O(1)$, a reduction by
factor $\approx 2/3$ relative to the unconstrained $6K$-dimensional space.
Icosahedral constraints with $N = 12$ reduce the space by factor $\approx
11/12$, leaving $O(K)$ free parameters at fixed order.

\subsubsection{Implementation: mode-selection projector}

Rather than assembling the Reynolds operator $P_G = |G|^{-1}\sum_{g\in G}\rho(g)$
at cost $O(|G|\,K^2)$, we exploit the explicit mode-selection rules derived in
MathFrame~\cite{MathFrame2026} Propositions~4.1--4.3.  For every point group
$G$ the fixed-point subspace $\Fix(G)$ is spanned by a subset of Fourier modes
that can be identified by a single modular-arithmetic lookup per mode index.
The resulting \texttt{build\_projector} routine runs in $O(K)$ and returns a
boolean mask rather than a dense matrix.

\begin{lstlisting}[style=python, caption={Mode-selection projector (MathFrame §4).}]
import jax.numpy as jnp
import numpy as np

# Mode-selection rules derived from MathFrame Propositions 4.1-4.3.
# Each entry: (xy_residue, z_residue, modulus)
# Mode k survives in xy iff k % modulus == xy_residue  (None = never)
# Mode k survives in z  iff k % modulus == z_residue   (None = never)
_MODE_RULES = {
    # Cyclic / pyramidal: 3D orbits possible
    "Cn":  lambda n: (1 % n, 0, n),
    "Cnv": lambda n: (1 % n, 0, n),
    # Dihedral / prismatic: z forced to zero (no 3D)
    "Dn":  lambda n: (1 % n, None, n),
    "Dnh": lambda n: (1 % n, None, n),
    # Cubic: sparse mode structure, 3D possible
    "T":   lambda _: (1, 0, 3),
    "Td":  lambda _: (1, 0, 3),
    "O":   lambda _: (1, 0, 4),
    "Oh":  lambda _: (1, 0, 4),
    "I":   lambda _: (1, 0, 5),
    "Ih":  lambda _: (1, 0, 5),
}

def build_projector(group: str, n: int, K: int):
    """
    Return boolean masks (xy_mask, z_mask) of shape (K,) indicating which
    Fourier modes k=1..K survive the symmetry constraint for the given group.

    Parameters
    ----------
    group : str   Schoenflies symbol, e.g. 'Cnv', 'T', 'Oh'
    n     : int   Fold number (ignored for cubic groups)
    K     : int   Truncation order

    Returns
    -------
    xy_mask : np.ndarray[bool], shape (K,)
    z_mask  : np.ndarray[bool], shape (K,)
    """
    rule = _MODE_RULES[group](n)
    xy_res, z_res, mod = rule
    ks = np.arange(1, K + 1)
    xy_mask = (ks % mod == xy_res) if xy_res is not None \
              else np.zeros(K, dtype=bool)
    z_mask  = (ks % mod == z_res)  if z_res  is not None \
              else np.zeros(K, dtype=bool)
    return xy_mask, z_mask

def pack_free(a_coeffs, b_coeffs, xy_mask, z_mask):
    """Extract the free (unconstrained) Fourier coefficients into a 1-D vector."""
    xy_a = a_coeffs[xy_mask, :2].ravel()   # x,y components of cos modes
    xy_b = b_coeffs[xy_mask, :2].ravel()   # x,y components of sin modes
    z_a  = a_coeffs[z_mask,  2]            # z component of cos modes
    z_b  = b_coeffs[z_mask,  2]            # z component of sin modes
    return np.concatenate([xy_a, xy_b, z_a, z_b])

def unpack_free(theta_free, K, xy_mask, z_mask):
    """Reconstruct full (K,2,3) coefficient array from free parameters."""
    n_xy = int(xy_mask.sum())
    n_z  = int(z_mask.sum())
    xy_a = theta_free[:2*n_xy].reshape(n_xy, 2)
    xy_b = theta_free[2*n_xy:4*n_xy].reshape(n_xy, 2)
    z_a  = theta_free[4*n_xy:4*n_xy+n_z]
    z_b  = theta_free[4*n_xy+n_z:]
    a = np.zeros((K, 3))
    b = np.zeros((K, 3))
    a[xy_mask, :2] = xy_a
    b[xy_mask, :2] = xy_b
    a[z_mask,   2] = z_a
    b[z_mask,   2] = z_b
    return jnp.array(np.stack([a, b], axis=1))   # (K, 2, 3)
\end{lstlisting}

The optimiser operates entirely on the \emph{free} parameter vector
$\theta_{\mathrm{free}}$ whose dimension is $\dim\Fix(G)$ (typically $O(K)$
with a group-dependent prefactor, as tabulated in §2.2.1).  Because the
parametrisation is intrinsic to $\Fix(G)$, no projection step is required
after gradient updates: every point in the free-parameter space corresponds
to a valid symmetry-constrained orbit.

\subsection{Action functional in JAX}

\begin{lstlisting}[style=python, caption={Core action functional (JAX).}]
import jax
import jax.numpy as jnp

def build_action(N: int, K: int, M: int, G: float = 1.0, m: float = 1.0):
    """
    Returns a JAX-jittable function action(theta) -> scalar.
    theta: shape (6*K,) -- packed Fourier coefficients [a1,b1,...,aK,bK],
           each a_k, b_k in R^3, centre-of-mass zeroed.
    """
    t = jnp.linspace(0, 1, M, endpoint=False)          # normalised time
    tau = jnp.arange(1, N) / N                          # phase offsets

    @jax.jit
    def action(theta):
        # Unpack coefficients: shape (K, 2, 3)
        coeff = theta.reshape(K, 2, 3)
        ks = jnp.arange(1, K + 1)[:, None]              # (K,1)

        # Synthesise gamma(t): shape (M, 3)
        cos_t = jnp.cos(2 * jnp.pi * ks * t[None, :])  # (K, M)
        sin_t = jnp.sin(2 * jnp.pi * ks * t[None, :])  # (K, M)
        gamma = jnp.einsum('ki,kc->ic', cos_t, coeff[:, 0, :]) \
              + jnp.einsum('ki,kc->ic', sin_t, coeff[:, 1, :])  # (M,3)

        # Kinetic term: (1/2) m |gamma_dot|^2
        gamma_dot = jnp.einsum('k,ki,kc->ic',
                               2*jnp.pi*ks[:, 0], -sin_t, coeff[:, 0, :]) \
                  + jnp.einsum('k,ki,kc->ic',
                               2*jnp.pi*ks[:, 0],  cos_t, coeff[:, 1, :])
        kinetic = 0.5 * m * jnp.mean(jnp.sum(gamma_dot**2, axis=-1))

        # Potential term: sum over phase offsets
        def potential_shift(s):
            cos_s = jnp.cos(2 * jnp.pi * ks * s)   # (K,)
            sin_s = jnp.sin(2 * jnp.pi * ks * s)
            gamma_s = jnp.einsum('k,ki,kc->ic', cos_s, cos_t, coeff[:,0,:]) \
                    + jnp.einsum('k,ki,kc->ic', sin_s, cos_t, coeff[:,0,:]) \
                    + jnp.einsum('k,ki,kc->ic', cos_s, sin_t, coeff[:,1,:]) \
                    + jnp.einsum('k,ki,kc->ic', sin_s, sin_t, coeff[:,1,:])
            diff = gamma - gamma_s                      # (M, 3)
            r = jnp.sqrt(jnp.sum(diff**2, axis=-1) + 1e-14)
            return -G * m**2 * jnp.mean(1.0 / r)

        potential = jax.vmap(potential_shift)(tau).sum()
        return kinetic + potential / N

    return action
\end{lstlisting}

\textbf{Note on collision regularisation.} The softening term $10^{-14}$ in
the denominator prevents gradient explosions during early optimisation but
must be removed before final Newton--Raphson refinement.  Orbits for which
$\min_t |\gamma(t) - \gamma(t+kT/N)| < 10^{-3}$ at convergence are flagged
as near-collision and discarded.

% ---------- 3. PHASE 1: ACTION MINIMISATION ----------
\section{Phase 1: Action Minimisation}

\subsection{Initialisation strategy}

\textbf{Algorithm: Batched symmetry-constrained initialisation.}
\begin{enumerate}[leftmargin=1.8em]
  \item[(In)] Point group $G$, body count $N$, batch size $B$, scale $\sigma$
  \item[(1)] Precompute Reynolds projector $P_G \in \R^{6K \times 6K}$
  \item[(2)] Sample $Z \sim \mathcal{N}(0, \sigma^2 I_{B \times 6K})$
  \item[(3)] $\Theta \leftarrow Z \cdot P_G^\top$ \hfill \textit{(project onto $\Fix(G)$)}
  \item[(4)] Normalise each row: $\Theta_i \leftarrow \sigma \Theta_i / \|\Theta_i\|$
  \item[(Out)] Batch $\Theta \in \R^{B \times 6K}$, all rows in $\Fix(G)$
\end{enumerate}

Default: $B = 8192$, $\sigma = 0.3$ (in units where $G = m = T = 1$).
The full batch is evaluated in a single \texttt{vmap}ped forward pass.

\subsection{L-BFGS via Optax}

Gradient descent uses the L-BFGS implementation in \texttt{optax} with
history length $h = 20$.  The gradient $\nabla_\theta \Action$ is provided
exactly by \texttt{jax.grad}.  After each step, the update is projected back
to $\Fix(G)$:
\[
  \theta \leftarrow P_G \bigl(\theta - \alpha \cdot \text{L-BFGS}(\nabla_\theta
  \Action)\bigr).
\]
Convergence is declared when $\|\nabla_\theta \Action\|_2 < 10^{-9}$.

\subsection{Candidate filtering}

After optimisation, candidates are filtered by:
\begin{enumerate}[leftmargin=1.8em]
  \item \textbf{Action threshold}: retain only candidates with
        $\Action < \Action_{\mathrm{planar}}(N)$, i.e.\ strictly below
        the best known planar choreography at the same $N$.
  \item \textbf{Planarity test}: compute the ratio of the smallest to largest
        principal moment of the point cloud $\{\gamma(t_i)\}$.  Discard if
        ratio $< 10^{-2}$.
  \item \textbf{Near-collision filter}: discard if any pairwise separation
        $< 10^{-3}$ along the refined orbit.
  \item \textbf{Duplicate detection}: cluster retained candidates by action
        value (tolerance $10^{-6}$) and symmetry group to remove re-discovered
        copies.
\end{enumerate}

% ---------- 4. PHASE 2: ORBIT REFINEMENT ----------
\section{Phase 2: Orbit Refinement}

\subsection{Yoshida 6th-order symplectic integrator}

Fourier-space minimisation yields an approximate periodic orbit.  It is
refined to a true solution of the $N$-body ODE using a Yoshida 6th-order
symplectic integrator \cite{Yoshida1990}, which preserves the symplectic
structure and gives energy conservation error $O(\Delta t^6)$.

The Yoshida coefficients $(w_0, w_1, \ldots, w_7)$ are fixed by the
root of a specific polynomial; we use the Yoshida (1990) set\,1 values,
stored as compile-time constants.  Step size $\Delta t = T/(200K)$.

\subsection{Newton--Raphson periodicity correction}

Let $\Phi_T : \R^{6N} \to \R^{6N}$ denote the time-$T$ flow map computed by
the symplectic integrator.  We seek $\mathbf{x}^* \in \R^{6N}$ such that
$\Phi_T(\mathbf{x}^*) = \mathbf{x}^*$.  Starting from the Fourier candidate,
Newton--Raphson iterations:
\[
  \mathbf{x}^{(n+1)} = \mathbf{x}^{(n)}
  - \left[D\Phi_T(\mathbf{x}^{(n)}) - I\right]^{-1}
    \left[\Phi_T(\mathbf{x}^{(n)}) - \mathbf{x}^{(n)}\right]
\]
are iterated until $\|\Phi_T(\mathbf{x}^*) - \mathbf{x}^*\|_\infty < 10^{-12}$.
The Jacobian $D\Phi_T$ is computed by integrating the variational equations
alongside the orbit (Section~\ref{sec:stability}).  Typically 3--5 Newton
steps suffice from a good Fourier initial point.

\subsection{Symmetry verification}

After refinement, the symmetry of the orbit is re-verified numerically:
for each generator $(R, \sigma) \in G$, compute
$\|R\,\bfq_1(t) - \bfq_1(t - \sigma T/N)\|_\infty$ and assert $< 10^{-9}$.
Orbits failing this check are re-classified under their actual (smaller)
symmetry group.

% ---------- 5. PHASE 3: KNOT CLASSIFICATION ----------
\section{Phase 3: Knot Classification Pipeline}

\subsection{Spacetime trace construction}

The spacetime trace $\Gamma$ is the closed curve in $\R^3$ obtained by
``unwinding'' the time axis:
\[
  \Gamma(\phi) = \bigl(\gamma(T\phi/2\pi),\;
                  R_z(\phi)\,\hat{e}_z\bigr)
                 \in \R^3,
  \quad \phi \in [0, 2\pi).
\]
In practice we embed $\Gamma$ in $S^3 \subset \R^4$ using the standard
Hopf-type map and project to $\R^3$ via stereographic projection before
passing to SnapPy.

\subsection{Projection and crossing diagram}

From the sampled curve $\Gamma$ (with $L = 2048$ sample points), we:
\begin{enumerate}[leftmargin=1.8em]
  \item Compute the writhe $w = \frac{1}{4\pi}\iint \frac{(\Gamma(s) -
        \Gamma(t))}{|\Gamma(s) - \Gamma(t)|^3} \cdot (d\Gamma(s) \times
        d\Gamma(t))$ numerically (discretised Gauss integral).
  \item Select a projection direction minimising the number of crossings
        (sampled over 200 random directions on $S^2$).
  \item Compute the Gauss code from the ordered crossing sequence and
        convert to Dowker--Thistlethwaite (DT) notation.
\end{enumerate}

\subsection{Polynomial invariants via SnapPy}

\begin{lstlisting}[style=python, caption={Knot invariant computation.}]
import snappy

def compute_knot_invariants(dt_notation: str) -> dict:
    """
    Given DT notation of the spacetime trace, compute invariants.
    """
    L = snappy.Link(dt_notation)
    K = L.components[0] if L.is_knot() else None

    results = {
        'alexander': str(L.alexander_polynomial()),
        'homfly_pt': str(L.homfly_polynomial()),
        'determinant': L.determinant(),
        'signature': L.signature(),
    }
    if K is not None:
        M = K.complement()
        results['hyperbolic_volume'] = M.volume() if M.solution_type() == \
            'all_positively_oriented_tetrahedra' else None
    return results
\end{lstlisting}

% ---------- 6. PHASE 4: STABILITY ----------
\section{Phase 4: Stability Analysis}
\label{sec:stability}

\subsection{Variational equations}

The linearised equations of motion for a displacement $\delta\mathbf{x}(t)$
from the reference orbit $\mathbf{x}^*(t)$ are:
\[
  \delta\dot{\mathbf{x}} = J(t)\,\delta\mathbf{x},
  \qquad
  J_{ij}(t) = \frac{\partial^2 H}{\partial x_i \partial x_j}\bigg|_{\mathbf{x}^*(t)},
\]
where $H$ is the $N$-body Hamiltonian.  These are integrated alongside the
orbit as a $6N \times 6N$ matrix ODE, yielding the fundamental matrix
$\Phi(t)$.  The monodromy matrix is $M = \Phi(T)$.

\begin{decision}
The Jacobian $J(t)$ is computed via \texttt{jax.hessian} of the Hamiltonian
at each timestep rather than by finite differencing, ensuring exact second
derivatives.
\end{decision}

\subsection{Floquet classification}

Eigenvalues of $M$ are classified as:
\begin{center}
\begin{tabular}{lll}
\toprule
\textbf{Type} & \textbf{Condition} & \textbf{Interpretation} \\
\midrule
Elliptic         & $|\lambda| = 1$, $\lambda \notin \R$  & Linearly stable mode \\
Hyperbolic       & $\lambda \in \R$, $|\lambda| \neq 1$  & Unstable mode \\
Complex-hyp.     & $|\lambda| \neq 1$, $\lambda \notin \R$ & Unstable, spiralling \\
Trivial          & $\lambda = 1$ (multiplicity $\geq 4$) & Symmetry-mandated \\
\bottomrule
\end{tabular}
\end{center}

An orbit is declared \emph{linearly stable} if all non-trivial eigenvalues
are elliptic.  The four trivial unit eigenvalues arise from time-translation
and spatial-translation symmetry.

\subsection{Lyapunov exponents}

For linearly unstable orbits, the leading Lyapunov exponent $\lambda_1$
is estimated by the standard QR-renormalisation method over $10^4$ periods.

% ---------- 7. VISUALISATION ----------
\section{Phase 5: Visualisation Pipeline}

\subsection{Tube rendering}

Orbital traces are rendered as smooth tubes in $\R^3$ using PyVista:
\begin{itemize}[leftmargin=1.8em]
  \item \textbf{Tube radius}: $0.02 \times$ characteristic orbit scale.
  \item \textbf{Colour map}: the tube is coloured by arclength parameter to
        make the body's progression visually legible.
  \item \textbf{Symmetry axes}: rotation axes are drawn as semi-transparent
        cylinders; mirror planes as translucent discs.
  \item \textbf{Knot diagram inset}: a 2-D knot diagram of the spacetime trace
        is composited as an inset in the lower-right corner of each render.
\end{itemize}

\subsection{Video export}

\begin{decision}
Output format: \texttt{H.264}, 3840$\times$2160 (4K), 60\,fps,
CRF\,18 (visually lossless). Container: \texttt{.mp4}.
\end{decision}

The animation shows the bodies moving along the pre-rendered tube, with a
camera orbit completing one revolution per two choreographic periods.
The \texttt{ffmpeg-python} wrapper handles frame assembly; PyVista's
off-screen renderer (\texttt{pyvista.start\_xvfb()}) is used in the
headless VM environment.

\subsection{Static publication figures}

Each orbit additionally produces a static 300\,dpi PDF figure (via PyVista's
\texttt{save\_graphic}) suitable for journal submission: three orthographic
views (top, front, side) plus the 3/4 perspective view, assembled on a
$2\times 2$ grid.

% ---------- 8. CATALOGUE FORMAT ----------
\section{Catalogue Format}

All discovered orbits are stored in a single HDF5 file \texttt{catalogue.h5}
with the following schema:

\begin{center}
\begin{tabular}{lll}
\toprule
\textbf{Dataset} & \textbf{Shape} & \textbf{Content} \\
\midrule
\texttt{orbits/\{id\}/q}         & $(M, N, 3)$  & Positions at $M$ timesteps \\
\texttt{orbits/\{id\}/p}         & $(M, N, 3)$  & Momenta \\
\texttt{orbits/\{id\}/action}    & scalar       & Action value \\
\texttt{orbits/\{id\}/period}    & scalar       & Period $T$ \\
\texttt{orbits/\{id\}/symgroup}  & string       & Sch\"onflies label \\
\texttt{orbits/\{id\}/knot}      & string       & Knot name or DT notation \\
\texttt{orbits/\{id\}/alexander} & string       & Alexander polynomial \\
\texttt{orbits/\{id\}/homfly}    & string       & HOMFLY-PT polynomial \\
\texttt{orbits/\{id\}/floquet}   & $(6N,)$      & Monodromy eigenvalues \\
\texttt{orbits/\{id\}/stable}    & bool         & Linear stability flag \\
\bottomrule
\end{tabular}
\end{center}

Initial conditions are stored to full IEEE\,754 double precision (64-bit).
The catalogue is accompanied by a \texttt{README.md} and a Python reader
module \texttt{choreo\_catalogue.py} for programmatic access.

% ---------- 9. MILESTONES ----------
\section{Development Milestones}

\begin{center}
\begin{tabular}{clc}
\toprule
\textbf{Milestone} & \textbf{Deliverable} & \textbf{Target} \\
\midrule
M0 & Environment setup, JAX GPU verify, test on figure-8 & Week 1 \\
M1 & $C_{3v}$ and $D_3$ search complete ($N=3$)           & Week 2 \\
M2 & All $N=3$ and $N=4$ point groups searched            & Week 4 \\
M3 & Knot pipeline validated; classification of M1/M2 results & Week 5 \\
M4 & Stability analysis complete for all discovered orbits & Week 6 \\
M5 & $N=6,12$ (tetrahedral, octahedral) search complete   & Week 9 \\
M6 & Icosahedral $N=12$ attempt                           & Week 11 \\
M7 & Visualisation pipeline, 4K renders, static figures   & Week 12 \\
M8 & Catalogue finalised, toolkit released, paper drafted & Week 14 \\
\bottomrule
\end{tabular}
\end{center}

% ---------- REFERENCES ----------
\begin{thebibliography}{9}

\bibitem{Yoshida1990}
H.~Yoshida,
\textit{Construction of higher order symplectic integrators},
Phys.\ Lett.\ A \textbf{150} (1990), 262--268.

\bibitem{CM2000}
A.~Chenciner and R.~Montgomery,
\textit{A remarkable periodic solution of the three-body problem in the case of equal masses},
Ann.\ Math.\ \textbf{152} (2000), 881--901.

\bibitem{Simo2002}
C.~Sim\'o,
\textit{New families of solutions in $N$-body problems},
in \textit{European Congress of Mathematics}, Birkh\"auser, Basel, 2001, pp.\ 101--115.

\bibitem{Palais1979}
R.~S.~Palais,
\textit{The principle of symmetric criticality},
Commun.\ Math.\ Phys.\ \textbf{69} (1979), 19--30.

\bibitem{MathFrame2026}
e5bed97b3349,
\textit{orbraid Mathematical Framework},
Technical Report v0.1, \url{https://github.com/e5bed97b3349/orbraid}, 2026.

\end{thebibliography}

\end{document}
